\documentclass{jstarticle}
\usepackage[utf8]{inputenc}
\usepackage[T1]{fontenc}
\usepackage[c3]{optidef}
\usepackage{lineno,hyperref}
\usepackage{xcolor}
\usepackage{algorithm}
\usepackage{comment}
\usepackage{amsmath}
\usepackage{tikz}
\usepackage{commath}
\usepackage{cases}
\usepackage{amsfonts}
\usepackage{amssymb}
\usepackage{algpseudocode}
\usepackage{subcaption}
\usepackage{multirow}
\usepackage{pdflscape}
\usepackage{array, booktabs, tabularx}
\usepackage{setspace}
\usepackage{booktabs}
\usepackage{rotating}
\usepackage{longtable}
\usepackage{float}
\usepackage{enumitem}
\usepackage{url}
\usepackage{diagbox}
\usepackage{makecell}
\usepackage{mathtools}
\usepackage{amsthm}
\theoremstyle{remark}
\newtheorem{remark}{Remark}
\newtheorem{theorem}{Theorem}
\newtheorem{assumption}{Assumption}
\addbibresource{refs.bib}
\usepackage{listings}
\lstset{
  basicstyle=\ttfamily\footnotesize,
  keywordstyle=\color{blue}\bfseries,
  commentstyle=\color{gray},
  stringstyle=\color{green!50!black},
  breaklines=true,
  frame=none,
  backgroundcolor=\color{gray!5},
  showstringspaces=false
}
\submitTo{ssad}

%% ---------------------------------------------------------------
%% Title
%% ---------------------------------------------------------------
\title{AuraBeam: An End-to-End Hardware-in-the-Loop Smart Matrix
       Headlight System Using Ensemble AI and Sensor Fusion}

%% ---------------------------------------------------------------
%% Authors
%% ---------------------------------------------------------------
\author{Nguyen Van Huan$^{1,*}$, Nguyen Thanh Kien$^{1}$, Duong Tan Minh$^{1}$,
         Dang Quang Huy$^{1}$, Vu Van Giang$^{1}$}

\affiliation{$^{1}$School of Information and Communication Technology,
             Hanoi University of Science and Technology, Hanoi, Vietnam}

\authornote{*}{Supervisor: Ho Viet Duc Luong, SOICT-HUST}

%% ---------------------------------------------------------------
%% Abstract
%% ---------------------------------------------------------------
\begin{abstract}
Nighttime road fatalities remain disproportionately high, with
approximately 50\% of fatal crashes occurring during periods that
account for only 25\% of total traffic volume.
A critical yet underaddressed failure mode in existing
Adaptive Driving Beam (ADB) systems is the
\emph{glare-induced detection loop failure}: when the front-facing
camera is saturated by oncoming headlights, the suppression bounding
box vanishes and the system involuntarily restores full high-beam
output, directly endangering the opposing driver.
This paper presents \textbf{AuraBeam}, an end-to-end
Hardware-in-the-Loop (HIL) simulation platform for smart matrix
headlight control that resolves this failure mode through three
core contributions.
\textbf{(C1)} An ensemble vision pipeline running YOLOv5 and YOLOv8
in parallel, fused via Weighted Non-Maximum Suppression, substantially
improves detection Recall under blooming and adverse illumination
conditions without a proportional increase in inference latency.
\textbf{(C2)} A six-state 3-D Kalman Filter fuses noisy pixel-space
coordinates $(x,y)$ from the vision pipeline with depth estimates $z$
from a physics-based Pseudo-Radar module. During camera occlusion, the
filter switches to a radar-only observation model and propagates the
suppression trajectory kinematically, eliminating LED flicker at zero
additional hardware cost.
\textbf{(C3)} The \emph{Glare Suppression Success Rate} (GSSR) metric
is introduced as a hardware-aware evaluation standard defined over the
$8{\times}8$ LED actuation grid via grid-level IoU, capturing
end-to-end actuation fidelity beyond what pixel-space mAP reflects.
Experiments on a five-scenario Golden Test Set covering urban night,
curved roads, rain, fog/snow, and lightning flash demonstrate that
AuraBeam achieves a mean GSSR of \textbf{93.1\%} and a trajectory
RMSE reduction of \textbf{68.6\%} relative to the YOLO-only baseline.
The fusion core executes in \textbf{0.61\,ms} on average, and the
complete pipeline is projectable to $>$30\,FPS on edge AI hardware.
Critically, in the lightning scenario where the YOLO-only baseline
collapses to a GSSR of 34.7\%, AuraBeam sustains \textbf{99.2\%}
GSSR, validating the indispensability of 3-D sensor fusion for
robust ADB control.
\end{abstract}

%% ---------------------------------------------------------------
%% Keywords
%% ---------------------------------------------------------------
\keyword{adaptive driving beam (ADB)}
\keyword{ensemble YOLO networks}
\keyword{Kalman filter, sensor fusion}
\keyword{hardware-in-the-loop simulation}
\keyword{glare suppression}
\keyword{LED matrix}
\keyword{advanced driver assistance systems (ADAS)}

\addbibresource{refs.bib}

%% ---------------------------------------------------------------
%% Extra packages
%% ---------------------------------------------------------------
\usepackage{tabularray}
\usepackage{xcolor}
\definecolor{customgray}{gray}{0.9}
\usepackage{tabularx}
\lstset{basicstyle=\ttfamily}
\let\Code\lstinline
\usepackage{prettyref}
\newrefformat{fig}{Fig.\,\ref{#1}}
\newrefformat{tab}{Table\,\ref{#1}}
\newrefformat{sec}{Section~\ref{#1}}
\newrefformat{subsec}{Section~\ref{#1}}
\newrefformat{subsubsec}{Section~\ref{#1}}

%% ---------------------------------------------------------------
\begin{document}
\maketitle

%% ==============================================================
\section{Introduction}
\label{sec:intro}
%% ==============================================================

Approximately 50\% of fatal road traffic accidents occur at night,
despite nighttime traffic accounting for only 25\% of total traffic
volume~\cite{fhwa2025nightvisibility}. Glare from oncoming headlights
has been identified as a primary contributing factor, degrading driver
reaction time by up to 1.4\,s~\cite{hwang2018headlightglare,
beard2025glarelitreview}. The automotive industry has addressed this
challenge through two generations of technology. \emph{Auto High Beam
(AHB)} systems toggle the full headlight beam based on overall scene
luminance, lacking spatial selectivity and introducing perceptible
flicker. More advanced \emph{Adaptive Driving Beam (ADB)} systems,
standardized under UN ECE Regulation 123, achieve pixel-level glare
suppression by selectively deactivating LED matrix segments directed
at opposing road users while preserving illumination in all other
regions~\cite{Shin2026DOE_AFS}.

Despite their efficacy, commercial ADB systems share a common
architectural vulnerability: they operate as purely reactive,
camera-centric feedback loops. When the front-facing camera is
directly saturated by an opposing headlight---a condition most
prevalent at close range, precisely when the safety stakes are
highest---the reliability of the vision algorithm collapses to near
zero. The suppression bounding box disappears, causing the system to
automatically restore full high-beam output and re-expose the opposing
driver to glare. This \emph{glare-induced detection loop failure} is a
fundamental fault mode that the majority of existing ADB algorithms
have not resolved. Furthermore, eliminating this failure without
relying on costly LiDAR or mmWave Radar sensor arrays remains a
significant barrier to democratizing this safety feature across
mainstream vehicle segments.

To overcome these limitations, this paper proposes \textbf{AuraBeam},
an end-to-end Hardware-in-the-Loop (HIL) simulation architecture for
smart matrix headlights, with three principal scientific contributions:

\begin{enumerate}[label=\textbf{C\arabic*.}]
    \item \textbf{Ensemble YOLO Vision Pipeline.}
          Two YOLOv5 and YOLOv8 models operate in parallel and their
          predictions are fused via Weighted Non-Maximum Suppression
          (Weighted NMS). This mechanism substantially improves
          detection Recall under adverse illumination conditions
          without a linear increase in processing latency.

    \item \textbf{3-D Kalman Filter Sensor Fusion.}
          A six-state Kalman Filter fuses noisy pixel-space coordinates
          $(x,y)$ from the vision pipeline with depth estimate $z$ from
          a physics-based Pseudo-Radar module. During occlusion, the
          filter switches to a radar-only update model and propagates
          the suppression trajectory kinematically, eliminating LED
          flicker entirely at zero additional hardware cost.

    \item \textbf{GSSR Metric and HIL Evaluation Framework.}
          The \emph{Glare Suppression Success Rate (GSSR)} metric,
          defined over an $8\times8$ LED grid using an IoU criterion,
          is introduced as a hardware-aware evaluation standard that
          more faithfully captures actuation fidelity than conventional
          mAP. A manually annotated Golden Test Set covering five
          adversarial scenarios is provided to ensure reproducibility.
\end{enumerate}

The remainder of this paper is organized as follows.
Section~\prettyref{sec:related} reviews related work.
Section~\prettyref{sec:method} details the AuraBeam methodology.
Section~\prettyref{sec:eval} defines the experimental evaluation
framework.
Section~\prettyref{sec:results} reports experimental results, and
Section~\prettyref{sec:conclusion} concludes.

%% ==============================================================
\section{Background and Related Work}
\label{sec:related}
%% ==============================================================

\subsection{Adaptive Driving Beam Systems}
\label{subsec:adb}

ADB technology was first commercialized in the 2013 Mercedes-Benz
S-Class under the designation ``Digital Light''~\cite{Bathla2022AutonomousVehicles}.
Under UN ECE R123, an ADB system is defined as a lighting installation
that provides a continuous high-beam while adapting to avoid glare for
other road users~\cite{Shin2026DOE_AFS}. Current implementations predominantly
employ Digital Micromirror Devices (DMDs) or segmented LED matrices at
considerable cost. Critically, all commercial systems are entirely
camera-centric; when input image quality degrades due to glare
saturation, the suppression mechanism fails
immediately~\cite{Kumar2015AutonomousVehicles}.

\subsection{Vehicle Detection and the Blooming Effect}
\label{subsec:detection}

In ADAS applications, the YOLO family of detection
models~\cite{Khanam2024YOLOv5, SCIRP126758} constitutes the industry
standard. However, nighttime headlight tracking is severely impeded by
the \emph{blooming effect}: at close range, the point-spread function of
an intense light source saturates neighboring pixels, completely
obliterating the geometric features surrounding the true bounding box.
To mitigate this degradation, ensemble detection strategies---running
multiple models in parallel and fusing their outputs---have been shown
to recover Recall under severe noise conditions, as demonstrated in UAV
tracking~\cite{Alshaer2025UAV} and pedestrian
detection~\cite{Dollar2014FastPyramids}.

\subsection{Trajectory Tracking and Sensor Fusion}
\label{subsec:kf}

Integrating a Kalman Filter (KF) into a detect-then-track paradigm is
an essential approach for robust object tracking~\cite{Kalman1960}.
Alshaer et al.~\cite{Alshaer2025UAV} demonstrated that coupling YOLO
with a KF reduces bounding box trajectory RMSE by up to 89\%.
Camera--Radar fusion has been widely adopted in collision avoidance and
lane-keeping
systems~\cite{ChavezGarcia2014SensorFusion}.

Applying a 3-D KF to the ADB problem, however, poses unique
challenges: \emph{(a)} when the camera is completely blinded, the
image-space state components become unobservable; \emph{(b)} depth
data must be optically independent to avoid correlated failure with the
vision channel. This paper proposes a physics-based Pseudo-Radar
module---operating independently of optical conditions---combined with
Kalman filter observability analysis, providing a viable approach for
cost-constrained ADB deployments that do not require additional
embedded sensors (real LiDAR/Radar).

\subsection{Hardware-in-the-Loop Testing}
\label{subsec:hil}

HIL simulation connects real embedded hardware with software plant
models, enabling algorithm validation prior to vehicle
integration~\cite{Kalman1960}. While HIL is standard practice
for brake ECU and steering system testing~\cite{Short2005HIL}, it is
rarely applied to matrix headlight control. Existing ADB evaluation
procedures rely primarily on laboratory photometric measurements or
on-road trials using expensive dedicated equipment. AuraBeam's HIL
framework, built on commodity microcontrollers and standard LED
matrices, offers a fully reproducible, cost-optimized alternative
for the research community.

%% ==============================================================
\section{System Methodology}
\label{sec:method}
%% ==============================================================

\subsection{System Overview}
\label{subsec:overview}

Fig.~\ref{fig:tong_quat} illustrates the AuraBeam processing pipeline.
Five modules operate within a tightly coupled feedback loop:

\begin{itemize}
    \item \textbf{Module A} --- Ensemble YOLO Vision Pipeline.
    \item \textbf{Module B} --- Pseudo-Radar Depth Estimation.
    \item \textbf{Module C} --- 3-D Kalman Filter Sensor Fusion.
    \item \textbf{Module D} --- Zone Mapping and Anti-Flicker Scheduling.
    \item \textbf{Module E} --- Arduino Hardware Actuation Layer.
\end{itemize}

\begin{figure*}[t] 
    \centering
    % Bạn có thể chỉnh 0.35 to nhỏ tùy ý (ví dụ 0.3 hoặc 0.4)
    \begin{subfigure}[b]{0.3\textwidth}
        \centering
        % [IMAGE REQUIRED: flow_charts/yolo_base.png]
        % This diagram shows the two parallel YOLO models (M1, M2) feeding into Weighted NMS
        \begin{center}
        \fbox{\textbf{[Figure 1: YOLO Ensemble Architecture - Image file required]}}
        \end{center}
        \caption{YOLO-only baseline pipeline.}
        \label{fig:anh_1}
    \end{subfigure}
    % --- ĐIỂM MẤU CHỐT NẰM Ở ĐÂY ---
    \hspace{2cm} % Thay \hfill bằng \hspace{2cm}. Tăng giảm số 2cm này để 2 ảnh xa/gần nhau theo ý thích!
    % -------------------------------
    \begin{subfigure}[b]{0.3\textwidth}
        \centering
        % [IMAGE REQUIRED: flow_charts/aurabeam.drawio.png]
        \fbox{\textbf{[Figure 2: AuraBeam System Pipeline - Image file required]}}
        \caption{AuraBeam full system pipeline.}
        \label{fig:anh_2}
    \end{subfigure}
    
    \caption{Architectural comparison between the YOLO-only baseline (left) and the full AuraBeam system (right).}
    \label{fig:tong_quat}
\end{figure*}

\begin{table}[H]
\centering
\caption{Hardware Bill of Materials}
\label{tab:bom}
\resizebox{\linewidth}{!}{
\begin{tabular}{llr}
\toprule
\textbf{Component} & \textbf{Specification} & \textbf{Est. Cost (USD)} \\
\midrule
Processing Unit    & PC / Laptop with discrete GPU & --- \\
Microcontroller    & Arduino Uno R3                & \$5 \\
LED Matrix         & 1588BS $8\times8$ (bare)      & \$1 \\
Camera             & USB 1080p Webcam / Dashcam    & \$15 \\
Serial Interface   & USB-A to USB-B cable          & \$2 \\
\midrule
\textbf{Total}     &                               & \textbf{\$23} \\
\bottomrule
\end{tabular}}
\end{table}

% ------------------------------------------------------------------
\subsection{Module A: Ensemble YOLO Vision Pipeline}
\label{subsec:yolo}
% ------------------------------------------------------------------

Two model variants, designated $\mathcal{M}_1$ (trained specifically
on high-glare headlight regions) and $\mathcal{M}_2$ (trained on a
general vehicle traffic dataset), run in parallel on each input
frame $I_t$. Each model produces a prediction set
$\mathcal{D}_i = \{(b_j, s_j) \mid j=1,\ldots,N_i\}$, where
$b_j = [c_x, c_y, w, h]^{\top}$ is the bounding box coordinate vector
and $s_j \in [0,1]$ is the associated confidence score.

The two prediction sets are merged and processed through
\emph{Weighted Non-Maximum Suppression}. For two overlapping boxes
$(b_a, s_a)$ and $(b_b, s_b)$ satisfying
$\mathrm{IoU}(b_a, b_b) \geq \tau_{\mathrm{NMS}}$, the fused box
$\hat{b}$ and score $\hat{s}$ are computed as:

\begin{equation}
\hat{b} = \frac{w_1 s_a b_a + w_2 s_b b_b}{w_1 s_a + w_2 s_b},
\label{eq:nms_box}
\end{equation}

\begin{equation}
\hat{s} = \frac{w_1 s_a + w_2 s_b}{w_1 + w_2},
\label{eq:nms_score}
\end{equation}

\noindent where $w_1, w_2 \in \mathbb{R}_{>0}$ are per-model weight
coefficients tuned on the validation set, and the IoU threshold is
fixed at $\tau_{\mathrm{NMS}} = 0.5$. The final output of Module A at
time $t$ is the image-space state vector:

\begin{equation}
\mathbf{d}_t = [c_x,\; c_y,\; w,\; h,\; \hat{s}]^{\top} \in \mathbb{R}^5.
\label{eq:det_vec}
\end{equation}

% ------------------------------------------------------------------
\subsubsection{Training Configuration}
\label{subsubsec:training}
% ------------------------------------------------------------------

The two constituent models are trained independently on purposefully
distinct data distributions to maximize feature complementarity.

\paragraph{Model $\mathcal{M}_1$ --- YOLOv5.}
$\mathcal{M}_1$ is trained on a custom, manually annotated dataset consisting of 5\,500 images. The training process is conducted with an input image size of $1024 \times 1024$ and a batch size of 16. The model is trained for 100 epochs, with an early stopping mechanism applied using a patience of 30 epochs to prevent overfitting. 

\paragraph{Model $\mathcal{M}_2$ --- YOLOv8.}
$\mathcal{M}_2$ is fine-tuned on another independent, in-house dataset comprising 6\,000 manually annotated images. Similar to the first model, $\mathcal{M}_2$ is trained for 100 epochs with an image resolution of 1024 and a batch size of 16. To significantly enhance model robustness and generalization, a comprehensive suite of data augmentation strategies is integrated during training. These include Mosaic (1.0), MixUp (0.2), Copy-Paste (0.1), and scaling (0.7). Additionally, color-space and spatial transformations are employed, specifically HSV saturation and value adjustments ($0.6$), random rotations ($5^{\circ}$), and perspective transformations ($0.0005$).

The ensemble weights $w_1, w_2$ in Eq.~\eqref{eq:nms_box} are tuned on the validation split to optimize the final detection performance, yielding the optimal weights for the combined predictions.
% ------------------------------------------------------------------
\subsection{Module B: Pseudo-Radar Depth Estimation}
\label{subsec:radar}
% ------------------------------------------------------------------

LiDAR and mmWave Radar sensors introduce significant cost and
integration complexity. Instead, AuraBeam deploys a
\emph{physics-based Pseudo-Radar} block that generates a depth signal
entirely independent of optical conditions. Assuming a Constant
Velocity (CV) approach model, the range estimate at time $t$ is:

\begin{equation}
Z(t) = Z_0 - v_{\mathrm{app}} \cdot t,
\label{eq:radar_z}
\end{equation}

\noindent where $Z_0$ is the initial range (estimated from the
bounding box size ratio at the first target acquisition) and
$v_{\mathrm{app}}$ is the relative approach velocity. Because $Z(t)$
is computed from clock cycles rather than optical signals, it is immune
to camera glare saturation.

\subsubsection*{Bounded Error Analysis of the Constant-Velocity Model}

The CV model has an inherent limitation: in practice, approaching
vehicles accelerate or brake. When the true acceleration
$a_{\mathrm{real}} \neq 0$, the true range evolves as:

\begin{equation}
Z_{\mathrm{true}}(t) = Z_0 - v_{\mathrm{app}} \, t
                       - \tfrac{1}{2}a_{\mathrm{real}} \, t^2,
\label{eq:radar_true}
\end{equation}

\noindent yielding a depth estimation error:

\begin{equation}
\Delta Z(t) = Z(t) - Z_{\mathrm{true}}(t)
            = -\tfrac{1}{2}a_{\mathrm{real}} \, t^2.
\label{eq:depth_error}
\end{equation}

Under moderate braking ($a \approx -0.5\,g$) over a 1\,s occlusion
window, the accumulated error is $|\Delta Z| \approx 2.45$\,m;
under emergency braking ($a \approx -0.9\,g$),
$|\Delta Z| \approx 4.4$\,m. These errors accumulate monotonically and
cannot be fully compensated by the Kalman filter during occlusion,
since the filter observes only the Pseudo-Radar signal (itself biased)
and the process model $\mathbf{Q}$ assumes CV. This constitutes a
\textbf{core limitation} of the current approach; future deployments
must adopt higher-order kinematic models or integrate real Radar
hardware.

% ------------------------------------------------------------------
\subsection{Module C: 3-D Kalman Filter Sensor Fusion}
\label{subsec:kf_fusion}
% ------------------------------------------------------------------

\subsubsection{State Representation and Covariance Matrices}

The Kalman filter maintains a six-dimensional state vector encoding
the object centroid coordinates and linear velocities along all three
axes:

\begin{equation}
\mathbf{x}_k = \bigl[c_x,\; \dot{c}_x,\; c_y,\; \dot{c}_y,\;
               Z,\; \dot{Z}\bigr]^{\top} \in \mathbb{R}^6.
\label{eq:state}
\end{equation}

The state transition matrix implements a constant-velocity model:

\begin{equation}
\mathbf{F} =
\begin{bmatrix}
1 & \Delta t & 0 & 0 & 0 & 0 \\
0 & 1 & 0 & 0 & 0 & 0 \\
0 & 0 & 1 & \Delta t & 0 & 0 \\
0 & 0 & 0 & 1 & 0 & 0 \\
0 & 0 & 0 & 0 & 1 & \Delta t \\
0 & 0 & 0 & 0 & 0 & 1
\end{bmatrix},
\label{eq:F_matrix}
\end{equation}

\noindent where $\Delta t = 33.3$\,ms is the frame interval. The process noise covariance matrix is initialized as a diagonal matrix, $\mathbf{Q} = \sigma_q^2 \mathbf{I}_6$ with $\sigma_q^2 = 0.01$ (empirically tuned to absorb small vehicle motion perturbations under the constant-velocity assumption). The measurement noise covariance matrices are:

\begin{align}
\mathbf{R}_{\mathrm{full}} &= \text{diag}(2.5, 2.5, 0.5) \text{ (px}^2\text{, px}^2\text{, m}^2\text{)}, \label{eq:R_full} \\
\mathbf{R}_{\mathrm{occ}} &= 10 \times \mathbf{R}_{\mathrm{full}},
\label{eq:R_occ}
\end{align}

\noindent reflecting high YOLO confidence during normal operation ($\sigma_x, \sigma_y = 2.5$ px, $\sigma_z = 0.5$ m) and degraded confidence during camera occlusion.

\subsubsection{Observation Models}

AuraBeam defines two observation matrices $\mathbf{H}$ and two
corresponding measurement noise matrices $\mathbf{R}$ for two
operating modes:

\noindent\textbf{Mode 1 --- Full Observation} (Camera and Radar
available):
\begin{equation}
\mathbf{H}_{\mathrm{full}} =
\begin{bmatrix}
1 & 0 & 0 & 0 & 0 & 0 \\
0 & 0 & 1 & 0 & 0 & 0 \\
0 & 0 & 0 & 0 & 1 & 0
\end{bmatrix}.
\label{eq:H_full}
\end{equation}
In this mode, $\mathbf{R}_{\mathrm{full}}$ is set low along the $X$
and $Y$ axes, reflecting high YOLO localization confidence.

\noindent\textbf{Mode 2 --- Occlusion} (Camera lost, Radar only):
\begin{equation}
\mathbf{H}_{\mathrm{occ}} =
\begin{bmatrix}
0 & 0 & 0 & 0 & 1 & 0
\end{bmatrix}.
\label{eq:H_occ}
\end{equation}

When $\hat{s}_t < \tau_{\mathrm{conf}}$, the system transitions to
Mode~2. Rather than completely removing the X--Y update equations,
an inflated measurement noise factor is injected
($\mathbf{R}_{\mathrm{occ}} = 10 \times \mathbf{R}_{\mathrm{full}}$).
This technique drives the Kalman Gain $\mathbf{K}$ toward zero in
image space, forcing the filter to rely exclusively on the internal
kinematic model and the Radar depth signal $Z$ to advance the
bounding box forward.

\subsubsection{Recursive Filter Equations}

\noindent\textbf{Prediction step (prior estimate):}
\begin{align}
\hat{\mathbf{x}}^{-}_{k} &= \mathbf{F}\,\hat{\mathbf{x}}_{k-1},
\label{eq:kf_pred_x} \\
\mathbf{P}^{-}_{k} &= \mathbf{F}\,\mathbf{P}_{k-1}\,\mathbf{F}^{\top}
                    + \mathbf{Q},
\label{eq:kf_pred_p}
\end{align}

\noindent where $\mathbf{Q}$ is the process noise covariance matrix.

\noindent\textbf{Update step (posterior estimate):}
\begin{align}
\mathbf{K}_k &= \mathbf{P}^{-}_k \mathbf{H}^{\top}
               \bigl(\mathbf{H}\mathbf{P}^{-}_k\mathbf{H}^{\top}
               + \mathbf{R}\bigr)^{-1},
\label{eq:kf_gain} \\
\hat{\mathbf{x}}_k &= \hat{\mathbf{x}}^{-}_k
                   + \mathbf{K}_k\bigl(\mathbf{y}_k
                   - \mathbf{H}\hat{\mathbf{x}}^{-}_k\bigr),
\label{eq:kf_update_x} \\
\mathbf{P}_k &= \bigl(\mathbf{I} - \mathbf{K}_k\mathbf{H}\bigr)
               \mathbf{P}^{-}_k,
\label{eq:kf_update_p}
\end{align}

\noindent where $\mathbf{H} \in
\{\mathbf{H}_{\mathrm{full}},\,\mathbf{H}_{\mathrm{occ}}\}$ is
selected by the mode-switching logic and $\mathbf{R}$ is the
corresponding measurement noise covariance.

\subsubsection{Occlusion Handling Algorithm}

Algorithm~\ref{alg:main} summarizes the complete
detect-then-track loop, including adaptive observation model switching.

\begin{algorithm}[H]
\caption{AuraBeam Detect-Track Loop with Adaptive Sensor Fusion}
\label{alg:main}
\begin{algorithmic}[1]
\Require Video stream $\{I_t\}$, thresholds
         $\tau_{\mathrm{conf}}$, $\tau_{\mathrm{NMS}}$,
         approach velocity $v_{\mathrm{ap}}$
\Ensure  LED $8\times8$ suppression grid $\mathbf{G}_t$ per frame
\State Initialize $\hat{\mathbf{x}}_0$, $\mathbf{P}_0$,
       $\mathbf{Q}$, $\mathbf{R}_{\mathrm{full}}$,
       $\mathbf{R}_{\mathrm{occ}}$
\State $Z_0 \gets$ estimate from initial bounding box size
\For{each frame $t = 1, 2, \ldots$}
    \State $\mathcal{D}_1, \mathcal{D}_2 \gets$ parallel YOLO
           inference on $I_t$
    \State $\mathbf{d}_t \gets$
           WeightedNMS$(\mathcal{D}_1 \cup \mathcal{D}_2,\;
           w_1, w_2, \tau_{\mathrm{NMS}})$
           \hfill\Comment{Eqs.~\eqref{eq:nms_box}--\eqref{eq:det_vec}}
    \State $Z(t) \gets Z_0 - v_{\mathrm{ap}} \cdot t$
           \hfill\Comment{Eq.~\eqref{eq:radar_z}}
    \State $\hat{\mathbf{x}}^-_t, \mathbf{P}^-_t \gets$
           KF\_Predict$(\hat{\mathbf{x}}_{t-1}, \mathbf{F},
           \mathbf{Q})$
           \hfill\Comment{Eqs.~\eqref{eq:kf_pred_x}--\eqref{eq:kf_pred_p}}
    \If{$\hat{s}_t \geq \tau_{\mathrm{conf}}$}
        \State $\mathbf{H} \gets \mathbf{H}_{\mathrm{full}}$;\quad
               $\mathbf{R} \gets \mathbf{R}_{\mathrm{full}}$
        \State $\mathbf{y}_t \gets [c_x,\; c_y,\; Z(t)]^{\top}$
    \Else
        \State $\mathbf{H} \gets \mathbf{H}_{\mathrm{occ}}$;\quad
               $\mathbf{R} \gets \mathbf{R}_{\mathrm{occ}}$
        \State $\mathbf{y}_t \gets [Z(t)]$
               \hfill\Comment{Radar-only propagation}
    \EndIf
    \State $\hat{\mathbf{x}}_t, \mathbf{P}_t \gets$
           KF\_Update$(\hat{\mathbf{x}}^-_t, \mathbf{K}_t,
           \mathbf{y}_t, \mathbf{H}, \mathbf{R})$
           \hfill\Comment{Eqs.~\eqref{eq:kf_gain}--\eqref{eq:kf_update_p}}
    \State $\mathbf{G}_t \gets$
           ZoneMap$(\hat{c}_x, \hat{c}_y, \hat{Z})$
           \hfill\Comment{Eq.~\eqref{eq:zone_map}}
    \State SerialSend$(\mathbf{G}_t)$
           \hfill\Comment{\texttt{BOX:c1:c2:r1:r2}}
\EndFor
\end{algorithmic}
\end{algorithm}

% ------------------------------------------------------------------
\subsection{Module D: Zone Mapping and Anti-Flicker Scheduling}
\label{subsec:zone}
% ------------------------------------------------------------------

\subsubsection{Pixel-to-Grid Coordinate Transform}

The Kalman-smoothed centroid $(\hat{c}_x,\,\hat{c}_y)$ is mapped to
an integer cell $(g_{\mathrm{col}},\,g_{\mathrm{row}})$ on the
$8\times8$ LED grid:

\begin{equation}
g_{\mathrm{col}} = \left\lfloor
  \frac{\hat{c}_x}{W_{\mathrm{frame}}} \cdot 8
\right\rfloor, \qquad
g_{\mathrm{row}} = \left\lfloor
  \frac{\hat{c}_y}{H_{\mathrm{frame}}} \cdot 8
\right\rfloor.
\label{eq:zone_map}
\end{equation}

A suppression region $\mathcal{B}$ with half-width $\delta = 1$ cell is
applied symmetrically around the centroid, creating a $3\times3$ cell suppression window. This conservative margin ensures complete coverage of the oncoming vehicle's light projection while maintaining reasonable host-driver forward visibility.

\subsubsection{Z-Axis Thresholding}

Table~\ref{tab:z_thresh} summarizes the three glare suppression modes
as a function of the estimated range $\hat{Z}$. These thresholds are
derived empirically from UN ECE R123 regulations: the 60\,m threshold
corresponds to the safe following distance for highway speeds (approximately
80\,km/h), while the 30\,m threshold represents the distance at which
opposing glare becomes critical for visual impairment according to
photometric standards.

\begin{table}[H]
\centering
\caption{Z-axis glare suppression activation thresholds}
\label{tab:z_thresh}
\resizebox{\linewidth}{!}{
\begin{tabular}{lll}
\toprule
\textbf{Estimated Range} & \textbf{LED State} & \textbf{Duty Cycle} \\
\midrule
$\hat{Z} \geq 60$\,m         & No suppression      & 100\% (full brightness) \\
$30 \leq \hat{Z} < 60$\,m    & Partial suppression & 40\% (dimming) \\
$\hat{Z} < 30$\,m            & Full suppression    & 0\% (selected zone off) \\
\bottomrule
\end{tabular}}
\end{table}

\subsubsection{Hysteresis Hold-Time}

To prevent \emph{ghosting}---rapid on/off toggling of LED cells at
zone boundaries---a hold-time condition is enforced:

\begin{equation}
\mathbf{G}_t[r,c] = \text{OFF}
\quad \text{if} \quad
\exists\, \tau \in [t-T_h,\,t]\!:\;
(r,c) \in \mathcal{B}_\tau,
\label{eq:hysteresis}
\end{equation}

\noindent where $T_h$ is the hold-time window, empirically set to
$T_h = 5$ frames (${\approx}167$\,ms at the evaluation frame rate).

% ------------------------------------------------------------------
\subsection{Module E: Hardware Actuation Layer}
\label{subsec:hw}
% ------------------------------------------------------------------

The Arduino Uno microcontroller scans the 1588BS LED matrix via
row-column multiplexing at 200\,Hz, well above the critical flicker
fusion threshold of the human visual system (60\,Hz). Control data
are transmitted from the host PC via a 115\,200-baud serial link
using the standardized packet format:

\begin{center}
\begin{verbatim}
BOX:<col_start>:<col_end>:
    <row_start>:<row_end>
\end{verbatim}
\end{center}

\noindent To prevent pipeline stalling, the firmware serial receiver
operates via an interrupt-driven ring buffer. Consequently, packet
parsing never preempts the timer ISR responsible for LED scanning,
guaranteeing stable light output regardless of upstream AI computation
latency.

%% ==============================================================
\section{Evaluation Framework and Metrics}
\label{sec:eval}
%% ==============================================================

\subsection{The GSSR Metric}
\label{subsec:gssr}

Conventional metrics such as mAP, Precision, and Recall characterize
detection accuracy in pixel space. For a matrix headlight controller,
however, the terminal output is a set of discrete LED cells: a small
pixel-space error may map to the wrong LED cell being suppressed,
resulting in residual glare or unnecessary darkening. A metric defined
directly over the LED grid is therefore required. The GSSR metric is
proposed to address this gap:

\begin{equation}
\mathrm{GSSR} = \frac{1}{N} \sum_{t=1}^{N}
\mathbf{1}\!\left[
  \mathrm{IoU}_{\mathrm{grid}}\!\left(\hat{\mathcal{B}}_t,\,
  \mathcal{B}^{*}_t\right) \geq 0.5
\right],
\label{eq:gssr}
\end{equation}

\noindent where $\hat{\mathcal{B}}_t$ is the set of LED cells predicted
by the system to require suppression at time $t$, and
$\mathcal{B}^{*}_t$ is the ground-truth set of cells corresponding to
the actual headlight footprint of the oncoming vehicle on the
$8 \times 8$ grid. Grid-level IoU is defined as:

\begin{equation}
\mathrm{IoU}_{\mathrm{grid}}(A,B) =
\frac{|A \cap B|}{|A \cup B|}.
\label{eq:grid_iou}
\end{equation}

The IoU threshold of 0.5 ensures that more than half of the glare
region is covered, satisfying the minimum criterion for glare reduction
specified in UN ECE R123 Article 6.2.

\subsection{Supplementary Metrics}
\label{subsec:metrics}

\noindent\textbf{False Darkening Rate (FDR):}
The proportion of frames in which the system suppresses at least one
LED cell when no opposing vehicle is present:

\begin{equation}
\mathrm{FDR} = \frac{1}{N_0}
\sum_{t:\, \mathcal{B}^{*}_t = \emptyset}
\mathbf{1}\!\left[\hat{\mathcal{B}}_t \neq \emptyset\right].
\label{eq:fdr}
\end{equation}

This metric quantifies unnecessary darkening, which degrades the host
driver's forward visibility.

\medskip
\noindent\textbf{Missed Glare Rate (MGR):}
The proportion of frames in which the system fails to suppress any LED
cell when an opposing vehicle is present:

\begin{equation}
\mathrm{MGR} = \frac{1}{N_1}
\sum_{t:\, \mathcal{B}^{*}_t \neq \emptyset}
\mathbf{1}\!\left[\hat{\mathcal{B}}_t = \emptyset\right].
\label{eq:mgr}
\end{equation}

MGR is the primary safety metric, representing events in which the
system fails to prevent glare exposure.

\medskip
\noindent\textbf{Trajectory RMSE:}
\begin{equation}
\mathrm{RMSE} = \sqrt{
  \frac{1}{N} \sum_{t=1}^{N}
  \left\|
    \begin{bmatrix}\hat{c}_{x,t} \\ \hat{c}_{y,t}\end{bmatrix}
    -
    \begin{bmatrix}c^{*}_{x,t} \\ c^{*}_{y,t}\end{bmatrix}
  \right\|_2^2
},
\label{eq:rmse}
\end{equation}

\noindent where $(c^{*}_{x,t}, c^{*}_{y,t})$ denotes the
ground-truth centroid coordinates.

\medskip
\noindent\textbf{Jitter Reduction (JR\%):}
\begin{equation}
\mathrm{JR} = \left(1 -
  \frac{\sigma_{\hat{c}}^2}{\sigma_{d}^2}
\right) \times 100\%,
\label{eq:jitter}
\end{equation}

\noindent where $\sigma_{\hat{c}}^2$ and $\sigma_{d}^2$ are the
centroid position variances of the Kalman Filter output and the raw
YOLO output, respectively.

\subsection{Golden Test Set Construction}
\label{subsec:testset}

A \emph{Golden Test Set} of 913 video frames is constructed via manual
annotation using Roboflow. Pixel-level bounding
boxes are subsequently projected onto LED grid cells via
Eq.~\eqref{eq:zone_map}. Table~\ref{tab:testset} summarizes the five
scenario groups.

Sequences are selected to represent adversarial edge cases that induce
failure in camera-only ADB systems:

\begin{itemize}
    \item \textbf{Urban Night.} Standard nighttime driving with
    unobstructed sightlines, serving as the control group for baseline
    detection performance.

    \item \textbf{Curved Roads.} Unlit mountain passes with successive
    bends and oncoming traffic, designed to stress the Kalman filter
    under nonlinear trajectories and evaluate lateral tracking error.

    \item \textbf{Rainy Night (Specular Reflection Noise).} Wet
    pavement generates specular reflections from opposing headlights,
    producing spurious light sources that challenge the vision pipeline's
    FDR suppression capability.

    \item \textbf{Fog/Snow (Scattering and Blooming).} Airborne
    particles scatter light and create diffuse bloom halos that
    obliterate vehicle geometric features. This scenario highlights the
    role of the ensemble, where $\mathcal{M}_1$ (light-source
    specialist) compensates for $\mathcal{M}_2$'s shape-feature
    degradation.

    \item \textbf{Lightning Flash (Exposure Disruption).} Sudden
    high-intensity illumination transients trigger rapid AEC adjustment,
    causing multi-frame detection dropout. This is the primary stress
    test for the 3-D Kalman filter's ability to sustain trajectory
    continuity via Pseudo-Radar during complete visual signal loss.
\end{itemize}

\begin{table}[H]
\centering
\caption{Golden Test Set scenario composition and environmental
         challenges}
\label{tab:testset}
\resizebox{\linewidth}{!}{%
    \begin{tabular}{llcl}
    \toprule
    \textbf{Scenario} & \textbf{Challenge Type}
      & \textbf{Frames} & \textbf{Primary Difficulty} \\
    \midrule
    Urban night      & Baseline                      & 204
      & Minimal environmental noise \\
    Curved road      & Nonlinear trajectory          & 78
      & Rapid lateral displacement \\
    Rainy night      & Specular reflection noise     & 219
      & Road reflections and lens water droplets \\
    Fog / Snow       & Strong light scattering       & 173
      & Blooming eliminates object features \\
    Lightning flash  & Abrupt illumination transient & 239
      & AEC disruption causes multi-frame dropout \\
    \midrule
    \textbf{Total}   &                               & \textbf{913} & \\
    \bottomrule
    \end{tabular}%
}
\end{table}

Fig.~\ref{fig:testset} illustrates representative frames from each
scenario, depicting the specular reflection, bloom scattering, and
exposure disruption conditions that constitute the primary failure
triggers for camera-only systems.

\begin{figure}[H]
\centering

\small\textbf{[Golden Test Set Representative Frames - Image files required]}

\medskip
Test set images required in \texttt{figs/} directory:
\begin{itemize}[leftmargin=*,label=--]
  \item (a) night\_frame.png: Urban night scenario with specular glare
  \item (b) curved\_frame.png: Curved road with oncoming vehicle
  \item (c) raining\_frame1.png: Heavy rain with water droplets
  \item (d) raining\_frame2.png: Rain continuation frame
  \item (e) snow\_frame.png: Fog/snow conditions with blooming
  \item (f) thunder\_frame1.png: Lightning flash with AEC disruption
  \item (g) thunder\_frame2.png: Lightning continuation frame
\end{itemize}

\caption{Representative Golden Test Set frames (image files pending).
(a)~Urban night; (b)~Curved road; (c--d)~Rainy night;
(e)~Fog/Snow; (f--g)~Lightning flash.}
\label{fig:testset}
\end{figure}

%% ==============================================================
\section{Experimental Results}
\label{sec:results}
%% ==============================================================

\subsection{Detection Performance (Module A)}
\label{subsec:res_det}

Table~\ref{tab:det_all} presents detection results for the three
model configurations---$\mathcal{M}_1$, $\mathcal{M}_2$, and
Ensemble---across all five Golden Test Set scenarios.

\begin{table}[H]
\centering
\caption{Detection performance comparison: single models vs.\ Ensemble
         across scenarios}
\label{tab:det_all}
\resizebox{\linewidth}{!}{%
\begin{tabular}{llccc}
\toprule
\textbf{Scenario} & \textbf{Configuration}
  & \textbf{mAP@0.5} & \textbf{Precision} & \textbf{Recall} \\
\midrule
\multirow{3}{*}{Urban night}
 & $\mathcal{M}_1$   & 53.12\% & 69.66\% & 52.57\% \\
 & $\mathcal{M}_2$   & 48.52\% & 68.82\% & 45.61\% \\
 & \textbf{Ensemble} & \textbf{48.74\%} & \textbf{53.57\%}
                     & \textbf{56.93\%} \\
\midrule
\multirow{3}{*}{Curved road}
 & $\mathcal{M}_1$   & 16.16\% & 55.32\% & 20.97\% \\
 & $\mathcal{M}_2$   & 45.08\% & 73.81\% & 50.00\% \\
 & \textbf{Ensemble} & \textbf{59.98\%} & \textbf{50.32\%}
                     & \textbf{62.90\%} \\
\midrule
\multirow{3}{*}{Rainy night}
 & $\mathcal{M}_1$   & 16.70\% & 31.97\% & 19.20\% \\
 & $\mathcal{M}_2$   & 32.42\% & 64.81\% & 27.62\% \\
 & \textbf{Ensemble} & \textbf{29.76\%} & \textbf{47.91\%}
                     & \textbf{38.82\%} \\
\midrule
\multirow{3}{*}{Fog / Snow}
 & $\mathcal{M}_1$   &  9.92\% & 12.61\% & 42.25\% \\
 & $\mathcal{M}_2$   & 18.14\% & 30.31\% & 28.01\% \\
 & \textbf{Ensemble} & \textbf{24.60\%} & \textbf{41.30\%}
                     & \textbf{44.88\%} \\
\midrule
\multirow{3}{*}{Lightning flash}
 & $\mathcal{M}_1$   & 39.76\% & 39.27\% & 54.88\% \\
 & $\mathcal{M}_2$   &  5.36\% & 13.85\% &  8.37\% \\
 & \textbf{Ensemble} & \textbf{49.51\%} & \textbf{31.45\%}
                     & \textbf{68.40\%} \\
\bottomrule
\end{tabular}%
}
\end{table}

\subsubsection*{The mAP--Recall Tradeoff: An IoU Drift Analysis}
\label{subsubsec:tradeoff}

A notable pattern in Table~\ref{tab:det_all} is that the Ensemble
mAP@0.5 does \emph{not} strictly dominate the best single model in
every scenario (e.g., Urban night: 48.74\% vs.\ 53.12\% for
$\mathcal{M}_1$). This outcome is a \emph{mathematically predictable
consequence} of Weighted NMS and is analyzed as follows.

When two overlapping boxes $(b_a, s_a)$ from $\mathcal{M}_1$ and
$(b_b, s_b)$ from $\mathcal{M}_2$ satisfy
$\mathrm{IoU}(b_a, b_b) \geq \tau_{\mathrm{NMS}}$, the fused box
$\hat{b}$ is the confidence-weighted centroid
(Eq.~\eqref{eq:nms_box}). Since $b_a \neq b_b$ in the general case,
$\hat{b}$ is displaced (drifted) from both source boxes and, by
extension, from the ground-truth annotation $b^*$. Formally, letting
$\Delta b = \hat{b} - b^*$ denote the residual displacement,
$\mathrm{IoU}(\hat{b}, b^*)$ decreases monotonically with
$\|\Delta b\|$. In practice, this IoU drift pushes borderline samples
below the 0.5 threshold, converting True Positives into False
Positives and thereby eroding mAP@0.5.

This tradeoff is, however, \textbf{by design}. The AuraBeam
architecture is governed by an ADAS fail-safe
philosophy~\cite{ISO26262_2018}:
\begin{itemize}[nosep]
  \item A \textbf{False Negative}---failing to detect an oncoming
    light source---causes the suppression zone to disappear, restoring
    full high-beam toward the opposing driver and creating a direct
    safety hazard.
  \item A \textbf{False Positive}---suppressing an LED cell when no
    vehicle is present---reduces the host driver's illumination
    marginally, without a safety consequence.
\end{itemize}
Under the principle of \emph{preferring over-suppression to
under-suppression}, running $\mathcal{M}_1$ and $\mathcal{M}_2$ in
parallel yields substantial Recall gains across adversarial scenarios:
from 8.37\% ($\mathcal{M}_2$ alone) to 68.40\% (Ensemble) in the
lightning scenario, and from 19.20\% to 38.82\% in the rainy night
scenario. In this context, \textbf{maximizing Recall carries a more
direct safety implication than optimizing mAP}, and the proposed GSSR
metric (Section~\ref{subsec:gssr}) is the appropriate evaluation
standard for LED-level actuation fidelity.

Four scenario-level observations follow. \textbf{(i)}~In Urban night
(baseline), $\mathcal{M}_1$ achieves peak mAP (53.12\%) under stable
illumination; the Ensemble improves Recall by 4.4 percentage points
(56.93\%) at an acceptable mAP cost. \textbf{(ii)}~In the Curved road
scenario, feature complementarity between the two models produces the
largest mAP gain ($+$14.9 points over $\mathcal{M}_2$), confirming
ensemble efficacy under nonlinear trajectories. \textbf{(iii)}~In the
Fog/Snow scenario, severe blooming reduces $\mathcal{M}_1$ Precision
to 12.61\%; the Ensemble recovers to 41.30\%, demonstrating
multi-model resilience. \textbf{(iv)}~The lightning scenario most
directly validates the dual-model architecture: $\mathcal{M}_2$
nearly collapses (mAP 5.36\%) while the Ensemble sustains 68.40\%
Recall, constituting direct empirical evidence that a single-model
architecture is insufficient for this problem.

\subsection{Tracking and Fusion Performance (Module C)}
\label{subsec:res_kf}

Table~\ref{tab:tracking_scenarios} compares trajectory RMSE and
jitter reduction across the three tracker configurations.
Fig.~\ref{fig:traj_lightning} and Fig.~\ref{fig:traj_rain} illustrate
representative trajectory plots for two characteristic occlusion
scenarios.

\begin{table}[H]
\centering
\caption{Tracking performance across the five Golden Test scenarios}
\label{tab:tracking_scenarios}
\resizebox{\linewidth}{!}{
\begin{tabular}{l ccc ccc}
\toprule
\multirow{2}{*}{\textbf{Scenario}}
  & \multicolumn{3}{c}{\textbf{RMSE (px) $\downarrow$}}
  & \multicolumn{3}{c}{\textbf{JR\% $\uparrow$}} \\
\cmidrule(lr){2-4} \cmidrule(lr){5-7}
  & YOLO & {+KF (2D)} & {AuraBeam}
  & YOLO & {+KF (2D)} & {AuraBeam} \\
\midrule
Urban night    & 36.8  & 16.5 & \textbf{16.5}
               & ---   & $-$3.4  & \textbf{$-$3.4} \\
Curved road    & 98.5  & 36.1 & \textbf{36.1}
               & ---   & 2.9     & \textbf{2.9} \\
Rainy night    & 238.1 & 80.2 & \textbf{80.2}
               & ---   & $-$10.5 & \textbf{$-$10.5} \\
Fog / Snow     & 86.9  & 25.1 & \textbf{25.1}
               & ---   & 12.7    & \textbf{12.7} \\
Lightning      & 78.7  & 11.7 & \textbf{11.7}
               & ---   & 16.3    & \textbf{16.3} \\
\midrule
\textbf{Mean}  & 107.8 & 33.9 & \textbf{33.9}
               & ---   & 3.6     & \textbf{3.6} \\
\bottomrule
\end{tabular}
}
\end{table}

\noindent\textbf{Note on the AuraBeam column.} The ``AuraBeam'' column
reflects the full configuration (Ensemble YOLO + 3-D KF +
Pseudo-Radar). In scenarios without complete occlusion (Urban night,
Curved road), the Pseudo-Radar contribution to planar $(x,y)$ RMSE
is negligible; its decisive role manifests along the Z-axis during
camera signal loss, as evidenced by the GSSR jump in
Table~\ref{tab:ablation}.

The 3-D Kalman filter reduces mean RMSE from 107.8\,px (YOLO only)
to 33.9\,px, a 68.6\% reduction. The highest JR\% values occur in the
lightning (16.3\%) and fog/snow (12.7\%) scenarios, confirming that
filter smoothing is most effective precisely when the vision signal is
most degraded. Negative JR\% in the urban night and rainy night
scenarios reflects a small phase lag introduced by the filter under
noise distributions that deviate from the assumed process model---an
acceptable tradeoff relative to the occlusion stability benefit.

\begin{figure}[H]
    \centering
    % [PDF REQUIRED: lightning_ablation_trajectory.pdf]
    \fbox{\textbf{[Trajectory Plot Required: lightning\_ablation\_trajectory.pdf]}}
    \caption{Centroid trajectory and RMSE analysis for the Lightning
             Flash scenario. The shaded region denotes the period of
             camera AEC disruption during which YOLO detection drops
             out entirely.}
    \label{fig:traj_lightning}
\end{figure}

\begin{figure}[H]
    \centering
    % [PDF REQUIRED: raining_ablation_trajectory.pdf]
    \fbox{\textbf{[Trajectory Plot Required: raining\_ablation\_trajectory.pdf]}}
    \caption{Centroid trajectory and RMSE analysis for the Heavy Rain
             scenario. AuraBeam sustains stable trajectory prediction
             across 88 consecutive frames of complete detection dropout.}
    \label{fig:traj_rain}
\end{figure}

\subsection{Occlusion Resilience: GSSR Results}
\label{subsec:res_gssr}

Table~\ref{tab:gssr} presents the primary result of this paper,
comparing the three system configurations across all five Golden Test
Set scenarios on the three actuation-level metrics: GSSR, FDR, and
MGR (defined in Section~\ref{subsec:metrics}).

\begin{table}[H]
\centering
\caption{GSSR, FDR, and MGR across the five test scenarios (\%)}
\label{tab:gssr}
\resizebox{\linewidth}{!}{%
\begin{tabular}{llccc}
\toprule
\textbf{Scenario} & \textbf{Configuration}
  & \textbf{GSSR}$\uparrow$ & \textbf{FDR}$\downarrow$
  & \textbf{MGR}$\downarrow$ \\
\midrule
\multirow{3}{*}{Urban night}
  & YOLO only         & 94.0\% &  4.9\% & 0.0\% \\
  & YOLO + KF (2D)    & 92.0\% &  7.4\% & 0.0\% \\
  & \textbf{AuraBeam} & \textbf{96.0\%} & \textbf{4.4\%}
                      & \textbf{0.0\%} \\
\midrule
\multirow{3}{*}{Curved road}
  & YOLO only         & 75.0\%  &  7.9\% & 3.9\% \\
  & YOLO + KF (2D)    & 80.3\%  &  3.9\% & 0.0\% \\
  & \textbf{AuraBeam} & \textbf{84.2\%} & \textbf{15.8\%}
                      & \textbf{0.0\%} \\
\midrule
\multirow{3}{*}{Rainy night}
  & YOLO only         & 46.1\%  & 12.8\% & 40.2\% \\
  & YOLO + KF (2D)    & 82.6\%  &  3.2\% &  3.2\% \\
  & \textbf{AuraBeam} & \textbf{87.2\%} & \textbf{9.6\%}
                      & \textbf{3.2\%} \\
\midrule
\multirow{3}{*}{Fog / Snow}
  & YOLO only         & 97.1\%  & 0.0\% & 0.0\% \\
  & YOLO + KF (2D)    & 100.0\% & 0.0\% & 0.0\% \\
  & \textbf{AuraBeam} & \textbf{98.8\%} & \textbf{1.2\%}
                      & \textbf{0.0\%} \\
\midrule
\multirow{3}{*}{Lightning flash}
  & YOLO only         & 34.7\%  &  0.4\% & 4.6\% \\
  & YOLO + KF (2D)    & 36.4\%  &  1.3\% & 0.0\% \\
  & \textbf{AuraBeam} & \textbf{99.2\%} & \textbf{0.8\%}
                      & \textbf{0.0\%} \\
\bottomrule
\end{tabular}%
}
\end{table}

Two principal findings emerge from Table~\ref{tab:gssr}.

\textbf{Breakthrough in the lightning scenario.} Under sudden
exposure disruption, the YOLO-only configuration degrades to a GSSR
of 34.7\%, and the YOLO + 2-D KF configuration improves only
marginally (36.4\%), as the 2-D filter lacks Z-axis information to
interpolate the trajectory kinematically during visual signal loss.
Integrating Pseudo-Radar depth elevates GSSR to \textbf{99.2\%}---an
absolute gain of $+$62.8 percentage points over both baselines---
confirming that 3-D tracking continuity through visual occlusion is
the decisive factor in the most dangerous ADB scenario.

\textbf{FDR tradeoff on curved roads.} AuraBeam achieves MGR = 0.0\%
but incurs an elevated FDR of 15.8\%, compared to 3.9\% for the 2-D
KF baseline. This results from the constant-velocity kinematic model's
linear prediction inertia: the filter maintains the bounding box at
the projected location even after the vehicle completes a curve.
Under the ADAS fail-safe principle---tolerating False Positives
(over-suppression) to prevent False Negatives (missed glare)---this
behavior is an intentional design property. A constant-turn-rate
kinematic model is identified as the primary mitigation in future
work.

Across the full Golden Test Set, AuraBeam achieves MGR = 0.0\% in 4
of 5 scenarios. The single exception (Rainy night, MGR = 3.2\%) still
represents a 37 percentage point reduction relative to the YOLO-only
baseline (40.2\%). These results confirm that the proposed 3-D sensor
fusion architecture satisfies the core safety requirement of the ADB
problem: \emph{no unmitigated glare exposure events} under adversarial
environmental conditions.

\subsection{Hardware Latency and Real-Time Performance}
\label{subsec:res_hw}

Table~\ref{tab:latency_all} presents the processing time budget across
the five scenarios. YOLO inference time includes preprocessing and
NMS post-processing for both parallel models, evaluated on a CPU-only
host in an offline setting.

\begin{table}[H]
\centering
\caption{Software latency and throughput analysis across five scenarios
         (offline evaluation environment)}
\label{tab:latency_all}
\resizebox{\linewidth}{!}{%
\begin{tabular}{lcccc}
\toprule
\textbf{Scenario}
  & \textbf{YOLO Inference (ms)}
  & \textbf{AuraBeam Fusion (ms)}
  & \textbf{Total Latency (ms)}
  & \textbf{FPS (offline)} \\
\midrule
Urban night    & 421.5 & 0.90 & 429.6 & 2.32 \\
Curved road    & 379.0 & 0.60 & 388.5 & 2.57 \\
Rainy night    & 396.6 & 0.50 & 403.2 & 2.48 \\
Fog / Snow     & 419.3 & 0.49 & 426.1 & 2.35 \\
Lightning      & 419.0 & 0.57 & 425.7 & 2.35 \\
\midrule
\textbf{Mean}  & \textbf{407.1} & \textbf{0.61}
               & \textbf{414.6} & \textbf{2.41} \\
\bottomrule
\end{tabular}%
}
\end{table}

The throughput bottleneck lies entirely in the unoptimized YOLO
inference block, which accounts for 98.2\% of total latency in the
offline environment. The AuraBeam fusion core (Modules B + C + D)
consumes a mean of \textbf{0.61\,ms}---0.15\% of the time
budget---with a standard deviation of $< 0.15$\,ms across all
weather conditions, confirming deterministic execution independent of
scene complexity.

For embedded deployment, it is necessary to account for the full
latency budget. Although the fusion kernel requires only 0.61\,ms, the
practical end-to-end latency is governed by YOLO inference
(${\sim}35$\,ms with TensorRT on Jetson Orin in optimal conditions)
and I/O communication. Serial UART at 115\,200 baud with a
${\sim}$30-byte packet yields a theoretical maximum of ${\sim}$14
messages/s (71\,ms/packet), constituting an additional bottleneck.
A realistic embedded deployment is projected to achieve 20--25\,FPS.
Meeting the ISO 26262 automotive real-time standard of
$\geq$30\,FPS~\cite{ISO26262_2018} requires: (1) TensorRT quantization,
(2) replacement of UART with CAN bus communication, and (3) latency
measurement on the actual embedded platform.

\subsection{Ablation Study}
\label{subsec:ablation}

Table~\ref{tab:ablation} quantifies the marginal contribution of each
module via a progressive component addition strategy, with GSSR and
JR\% averaged over the full Golden Test Set.

\begin{table}[H]
\centering
\caption{Ablation study: marginal contribution of each module
         (mean over the full Golden Test Set)}
\label{tab:ablation}
\resizebox{\linewidth}{!}{%
\begin{tabular}{lccc}
\toprule
\textbf{System Configuration}
  & \textbf{GSSR} & \textbf{$\Delta$GSSR} & \textbf{JR\%} \\
\midrule
(1) YOLO only (Baseline)          & 69.4\%  & ---            & 0.0\% \\
(2) $+$ 2-D Kalman Filter         & 78.3\%  & $+$8.9\%       & 3.6\% \\
\textbf{(3) $+$ Pseudo-Radar (AuraBeam)}
  & \textbf{93.1\%} & \textbf{$+$14.8\%} & \textbf{3.6\%} \\
\bottomrule
\end{tabular}%
}
\end{table}

The ablation data reveal two distinct contribution levels.
Step~(1)$\to$(2): the 2-D Kalman filter improves geometric stability
in image space, yielding a GSSR gain of 8.9 percentage points and a
jitter reduction of JR\% = 3.6\%.
Step~(2)$\to$(3): Pseudo-Radar integration is the decisive improvement,
delivering the largest single GSSR gain ($+$14.8 points) without
altering JR\%---confirming that Pseudo-Radar does not contribute to
2-D trajectory smoothing but instead resolves an independent problem:
maintaining Z-axis continuity of the suppression zone during visual
occlusion. This result supports the central thesis that the ADB
problem requires a full 3-D state representation, not merely a 2-D
smoothing of the projected vision signal.

%% ==============================================================
\section{Conclusion}
\label{sec:conclusion}
%% ==============================================================

This paper presented \textbf{AuraBeam}, an end-to-end
Hardware-in-the-Loop smart matrix headlight system that resolves the
glare-induced detection loop failure inherent in all camera-centric ADB
architectures. Three contributions were made: \textbf{(i)} an Ensemble
YOLOv5/v8 pipeline with Weighted NMS that improves detection Recall
under blooming conditions; \textbf{(ii)} a 3-D Kalman Filter with
adaptive dual observation models fusing vision coordinates with
Pseudo-Radar depth, enabling trajectory continuation through complete
visual occlusion; and \textbf{(iii)} the GSSR metric, a
hardware-aware evaluation standard that captures actuation fidelity at
the LED grid level. AuraBeam achieves a GSSR of 93.1\%, a trajectory
RMSE of 33.9\,px, and a fusion core latency of 0.61\,ms on a \$23
hardware platform. In the lightning scenario---the most dangerous
failure mode---AuraBeam sustains 99.2\% GSSR while the YOLO-only
baseline collapses to 34.7\%, validating the indispensability of 3-D
sensor fusion for robust ADB control.

\paragraph{Limitations and Safety Considerations.}
\begin{itemize}[leftmargin=*]
  \item \textbf{Constant-Velocity Model.} The Pseudo-Radar assumes
    constant $v_{\mathrm{app}}$. Under moderate braking
    ($a = -0.5\,g$), depth estimation error accumulates as
    $\Delta Z(t) \approx \frac{1}{2}|a|t^2$, reaching ${\sim}$1--2\,m
    after 1--1.5\,s of camera occlusion. Future work must implement
    adaptive velocity estimation or higher-order kinematic models
    (e.g., constant-turn-rate). A formal safety analysis linking
    maximum tolerable $\Delta Z$ to ADB control margins has not been
    performed.
  \item \textbf{$8\times8$ LED Resolution.} The prototype operates
    with 64 segments ($11.25°$ per segment), substantially coarser
    than production standards ($\geq1°$ per SAE J3014). This directly
    inflates FDR; scaling to a production-grade matrix geometry is a
    prerequisite for regulatory acceptance.
  \item \textbf{Single-Vehicle Tracking.} The Kalman filter assumes
    a single oncoming vehicle. Multi-vehicle scenarios (highway
    overtaking) require independent KF tracks per target, increasing
    computational load and synchronization complexity.
  \item \textbf{GSSR Metric Validation.} The proposed GSSR (grid-level
    IoU $\geq 0.5$) captures grid coverage but has \emph{not been
    validated} against photometric standards (illuminance in lux per
    UN ECE R123 Article 6). Bridging GSSR to lux compliance requires
    laboratory photometric measurement---a critical gap for genuine
    safety claims.
\end{itemize}

\paragraph{Future Work and Production Pathway.}
\begin{enumerate}[label=\textbf{(\roman*)},leftmargin=*]
  \item \textbf{Advanced Motion Models.} Replace the constant-velocity
    assumption with constant-turn-rate (CTR) or Wiener process
    acceleration models to handle real-world vehicle maneuvers.
    Quantify GSSR on curved-road scenarios where FDR currently reaches
    15.8\%.
  \item \textbf{Real Sensor Fusion.} Integrate a 77--79\,GHz mmWave
    Radar in place of the Pseudo-Radar simulation; benchmark depth
    accuracy against LiDAR ground truth, targeting $|\Delta Z| < 0.5$\,m
    under emergency braking.
  \item \textbf{Full Embedded Deployment.} Port to NVIDIA Jetson Orin
    Nano with TensorRT optimization; measure end-to-end latency from
    camera capture to LED actuation under real-time OS constraints.
    Validate $\geq$30\,FPS compliance per ISO 26262~\cite{iso26262}.
  \item \textbf{Production LED Hardware.} Scale from the $8\times8$
    prototype (64 segments) to a production-grade matrix
    ($\geq$512 segments); re-evaluate GSSR and FDR scaling laws under
    finer spatial granularity.
  \item \textbf{Photometric Validation.} Establish a goniophotometric
    bench with calibrated luminance meters; correlate the proposed GSSR
    metric with actual illuminance reduction (lux) per UN ECE R123
    Article 6.2. This step is mandatory for regulatory compliance claims.
  \item \textbf{Functional Safety Analysis.} Conduct ISO 26262
    Functional Safety Assessment (ASIL rating, FMEA). Characterize
    failure modes (Pseudo-Radar corruption, Kalman filter divergence)
    and implement fail-safe fallback strategies (e.g., conservative
    full-beam cutoff).
\end{enumerate}

%% ==============================================================
\printbibliography
%% ==============================================================

\end{document}